%% ====================================================================
%%  SECTION 10 --- RELATED WORK (~2pp)
%% ====================================================================
\section{Related Work}
\label{sec:related}

\paragraph{Where does the time-centric viewpoint come from?}
\citet{achille2025agents} establish that the mutual information between an agent's history and its decisions bounds expected speedup over random search.
Their framework is model-independent: it characterizes a task's information structure without reference to which model processes it.
The older lineage runs through Solomonoff induction, Levin search, and Kelly's~\citeyearpar{kelly1956new} result that mutual information quantifies the value of a noisy side channel for log-optimal betting~\citep{cover2006elements}---the idea that prior weight and reusable structure trade against search time.
We inherit the time-centric transductive view almost entirely.
Our contribution relative to Achille--Soatto is narrow and operational: they diagnose how much structure a task contains; we tell a builder which system to build to exploit that structure.
The moment one asks ``which model should I use?'', one needs per-step cost and within-mode competence---terms their framework does not contain.

\paragraph{How do compound AI systems handle routing today?}
The compound AI systems movement~\citep{zaharia2024compound,wen2025compound} recognizes that production systems combine multiple models, tools, and orchestration layers.
Within this space, a growing literature learns routing policies end-to-end: RouteLLM trains a classifier to select between models using preference data~\citep{ong2024routellm}; FrugalGPT cascades from cheap to expensive models~\citep{chen2023frugalgpt}; xRouter uses reinforcement learning for cost-aware orchestration~\citep{xrouter2024}; and OmniRouter applies Lagrangian decomposition to the routing \emph{assignment} problem~\citep{omnirouter2025}.
Cost-aware contrastive routing~\citep{cscr2025} and multi-agent routing~\citep{masrouter2025} extend these ideas to richer settings.

OmniRouter is the closest point of contact with our Lagrangian non-connection: it applies Lagrangian decomposition to the routing \emph{assignment} sub-problem, but decomposes the routing sub-problem, not the performance-gap identity.
None of these works provides a diagnostic decomposition of \emph{why} one design beats another.
They solve the routing problem algorithmically; we decompose the objective analytically.
The approaches are complementary: our decomposition identifies which lever matters, and their learning algorithms pull it.
The routing-triviality result (Section~\ref{sec:routing-triviality}) provides a pre-screen these methods lack: check whether routing adds value before investing in learning a router.
The gap is methodological---the routing literature treats model selection as a learning problem; this paper treats it as a measurement problem.

\paragraph{What do test-time compute scaling results tell us?}
A parallel literature demonstrates that smaller models with more inference-time compute can match or outperform larger models.
\citet{snell2025testtime} show that optimally allocating test-time compute can be more effective than scaling model parameters.
\citet{wu2025inference} derive empirical inference scaling laws.
\citet{brown2024monkeys}---in work titled \emph{Large Language Monkeys}---provide direct empirical support for the geometric-trials assumption underlying our graph-depth model: repeated independent sampling with verification produces pass rates well-described by $1-(1-p)^D$.
The $4/\delta$ bound~\citep{fourdelta2024} and OSCA~\citep{osca2024} extend this to sample-efficient allocation strategies.

These papers establish the empirical phenomenon that drives our depth-cost tradeoff.
None provides a closed-form characterization of optimal depth $\Dstar(M)$ or proves that optimal depth decreases monotonically with model competence.
Their approach is empirical compute-optimal allocation; ours is analytical.

\paragraph{What can the algorithm-selection literature contribute?}
The question ``how should one allocate limited compute across competing strategies?'' has a long OR history.
Weitzman's Pandora's Box~\citep{kaufmann2016bestarm} gives optimal stopping rules for inspecting alternatives with known reward distributions---structurally close to our routing problem, but static: there is no learning-from-verification feedback.
Algorithm portfolios and per-instance selection~\citep{multiagentdesign2025,qiao2025scaling} allocate resources across solvers, and racing methods~\citep{kaufmann2016bestarm} sequentially eliminate inferior candidates.
We bring these ideas into the language of modern agent systems, where the ``strategies'' are models with different cost-competence profiles and the ``budget'' is a deployment compute ceiling.
The classical results assume known reward distributions.
We add a sample-complexity layer they do not address, and the Bernstein corollary makes that layer practically feasible.

\paragraph{How do scaling laws and benchmarks relate?}
The scaling-law literature~\citep{kaplan2020scaling,hoffmann2022chinchilla} provides the cost-capability relationship that our graph-depth model takes as input: model cost and competence as functions of scale.
On the evaluation side, benchmark critiques~\citep{linegoes2025} document how aggregate metrics conflate distinct performance factors---precisely the conflation our decomposition resolves.
Lessons from LLM routing~\citep{chen2025optimizing} highlight that model selection depends on task distribution in ways that benchmark averages cannot capture.

Scaling laws characterize individual models.
They do not characterize systems.
The gap between ``how does a model scale?'' and ``how does an agent system scale?'' is the gap this paper addresses.

\paragraph{What does transfer learning theory say about reusing agent designs?}
Domain adaptation~\citep{bendavid2010theory} provides the gold-standard formalism for bounding target-domain risk in terms of source-domain risk plus a divergence penalty.
Transferability estimation using Bhattacharyya class separability~\citep{pandy2022gbc} brings these ideas to practical image-classification pipelines.
Our transfer trichotomy (Section~\ref{sec:transfer}) operates on different mathematical objects---design tuples and routing policies rather than hypothesis classes and feature representations---but the conceptual structure is parallel: two domains are close when their task-relevant structure is close.
The qualitative framework we present does not compete with the formal bounds of \citet{bendavid2010theory}; it addresses a different question (when does an \emph{agent system design} transfer?) that their theory does not cover.

\paragraph{Broader context.}
Agent design as a budgeted optimization problem has roots in stochastic programming, team decision theory~\citep{multiagentfinance2026}, and the Dantzig--Wolfe decomposition tradition~\citep{dantzigwolfe1960,fisher2004lagrangian}---though as Section~\ref{sec:lagrangian} shows, the analogy to Lagrangian dual decomposition is precisely wrong for our four-term identity.
